\clearpage
\VerdierPhysicalPageAnchor{0115}%
\begingroup
\small
\setlength{\abovedisplayskip}{4pt plus 2pt minus 2pt}%
\setlength{\belowdisplayskip}{4pt plus 2pt minus 2pt}%

\VerdierSemanticAnchor{ch2-1-2-3-proof-f1-f2}
具体的に、$f_1\colon X\longrightarrow X'$ と
$f_2\colon Y\longrightarrow Y'$ が同型である場合を考える。$\mathcal D$ の
任意の対象 $W$ に対し、次の可換図式がある。

\VerdierSemanticAnchor{ch2-1-2-3-five-lemma-diagram}
\begin{equation*}
  \resizebox{0.98\linewidth}{!}{%
    $\begin{tikzcd}[ampersand replacement=\&,column sep=2.15em,row sep=3em]
      \cdots \arrow[r]
        \& \operatorname{Hom}_{\mathcal D}(W,X) \arrow[r] \arrow[d,"{\operatorname{Hom}_{\mathcal D}(W,f_1)}"']
        \& \operatorname{Hom}_{\mathcal D}(W,Y) \arrow[r] \arrow[d,"{\operatorname{Hom}_{\mathcal D}(W,f_2)}"']
        \& \operatorname{Hom}_{\mathcal D}(W,Z) \arrow[r] \arrow[d,"{\operatorname{Hom}_{\mathcal D}(W,f_3)}"']
        \& \operatorname{Hom}_{\mathcal D}(W,X[1]) \arrow[r] \arrow[d,"{\operatorname{Hom}_{\mathcal D}(W,f_1[1])}"']
        \& \cdots \\
      \cdots \arrow[r]
        \& \operatorname{Hom}_{\mathcal D}(W,X') \arrow[r]
        \& \operatorname{Hom}_{\mathcal D}(W,Y') \arrow[r]
        \& \operatorname{Hom}_{\mathcal D}(W,Z') \arrow[r]
        \& \operatorname{Hom}_{\mathcal D}(W,X'[1]) \arrow[r]
        \& \cdots
    \end{tikzcd}$%
  }\,.
\end{equation*}

\VerdierSemanticAnchor{ch2-1-2-3-five-lemma}
各行は完全である。五項補題により、
$\operatorname{Hom}_{\mathcal D}(W,f_3)$ は任意の $W$ に対して同型であり、
したがって $f_3$ も同型である。系の残る二つの主張は、公理 TRII を用いて
この場合から従う。

\VerdierSemanticAnchor{ch2-1-2-4}
\medskip
\noindent\textbf{系 1.2.4.}\quad
$X\xrightarrow{u}Y\xrightarrow{v}Z\xrightarrow{w}X[1]$ と
$X\xrightarrow{u}Y\xrightarrow{v'}Z'\xrightarrow{w'}X[1]$ を二つの
完全三角とする。このとき、次の図式を可換にする同型
$f\colon Z\xrightarrow{\sim}Z'$ が存在する。

\VerdierSemanticAnchor{ch2-1-2-4-1}
\begin{equation}
  \begin{tikzcd}[column sep=2.8em,row sep=1.7em]
    && Z \arrow[dr,"w"] \arrow[dd,"f"',"\sim"] & \\
    X \arrow[r,"u"] & Y \arrow[ur,"v"] \arrow[dr,"v'"'] && X[1] \\
    && Z' \arrow[ur,"w'"'] &
  \end{tikzcd}
  \tag{1.2.4.1}
\end{equation}

\VerdierSemanticAnchor{ch2-1-2-4-proof}
実際、公理 TRIII により、図式 (1.2.4.1) を可換にする射
$f\colon Z\longrightarrow Z'$ が存在する。(1.2.3) により、この射は同型で
ある。

\VerdierSemanticAnchor{ch2-1-2-4-essential-uniqueness}
したがって系 (1.2.4) は、与えられた射を含む完全三角が本質的に一意、
すなわち同型を除いて定まることを示す。ただし、この同型自体は一意では
ない。この点は (1.2.12) で明確にする。

\VerdierSemanticAnchor{ch2-1-2-5}
\medskip
\noindent\textbf{系 1.2.5.}\quad
二つの完全三角の直和は完全三角である。

\VerdierSemanticAnchor{ch2-1-2-5-proof-start}
次の二つをとる。

\VerdierSemanticAnchor{ch2-1-2-5-input-triangles}
\[
  \begin{gathered}
    X\xrightarrow{u}Y\xrightarrow{v}Z\xrightarrow{w}X[1],\\
    X'\xrightarrow{u'}Y'\xrightarrow{v'}Z'\xrightarrow{w'}X'[1]
  \end{gathered}
\]
\endgroup
